\abstractheading
  {Defying Scaling in Potential Theory}
  {Iwona Chlebicka}
  {\affilUniversityWarsaw}
  {Invited lecture}

Classical non-linear potential theory thrives under the predictable comfort of
power-law scaling, where $p$-Laplacian operators yield well-behaved pointwise
estimates through traditional Wolff potentials. But what happens when homogeneity
breaks down completely? In this talk, we defy the constraints of standard scaling
to tackle measure data elliptic problems governed by second-order operators with
non-standard Orlicz growth and merely measurable coefficients.

We break away from homogeneous paradigms by establishing sharp pointwise
estimates expressed in terms of novel, generalized potentials tailored to
non-scaling structures. Going beyond describing local solution behavior where
standard tools fail, we uncover the exact regularity consequences dictated by
these estimates. Crucially, this framework holds across both scalar equations
and vectorial systems, forging a powerful potential theory liberated from
scaling laws.
